% !TeX root = ../main.tex
% Add the above to each chapter to make compiling the PDF easier in some editors.

\chapter{Introduction}\label{chapter:introduction}

\section{Background}

Social media platforms have emerged as primary channels for the dissemination of vast quantities of multimedia
content, including images, text, and video. These digital streams provide a rich source of data for research in the
social sciences, particularly in fields such as political extremism, regional conflict analysis, and hate speech
detection. To facilitate such research, scholars often rely on web crawlers to harvest large datasets from online
platforms. However, the raw data collected through these methods are typically unstructured and lack semantic
annotation, requiring significant preprocessing to extract meaningful insights.

For image data in particular, classification serves as a critical preprocessing step, transforming unstructured
content into labeled datasets that are suitable for both quantitative and qualitative analysis. Among the various
categories of harmful or sensitive visual content, Nazi-related imagery remains especially salient. Symbols such as
the swastika, once co-opted by the German National Socialist regime (1933–1945), continue to function as enduring
emblems of hate and extremism. Additional symbols—such as the \ac{SS} bolts, the Black Sun (linked to fascist occult
ideologies), and the Siegrune—have been appropriated by modern neo-Nazi and far-right movements. These symbols
persist across digital platforms, where they are employed to signal ideological affiliation, propagate extremist
views, or incite violence. Given that many countries restrict or criminalize the display of such imagery, there is a
pressing need for automated detection tools to support content moderation and policy enforcement.

\section{Problem Statement}

The automatic detection and classification of Nazi symbols in images is critical for researchers, journalists, civil
society organizations, and online platforms working to monitor and mitigate the spread of extremist content. While
commercial image classification services exist, they are often proprietary, lack transparency, and entail high usage
costs—particularly for applications requiring large-scale image analysis or custom symbol detection. This presents a
significant barrier not for governments per se, but for non-state actors and research groups who often operate under
tight computational and financial constraints.
Even when national institutions have greater resources, scalability and cost-efficiency remain key considerations.
High-throughput image classification tasks, such as analyzing large volumes of online content or social media images
in near real-time, can demand substantial compute time and infrastructure. In practice, this limits the feasibility
of using closed, black-box services that are either too slow, too expensive to scale, or unsuitable for
customization. Moreover, reliance on opaque systems may hinder accountability and reproducibility—both essential in
research and regulatory contexts.
To address these challenges, there is a need for an open, transparent, and efficient image classification system
tailored specifically for Nazi symbol detection. Such a system should support both binary classification: detecting
presence/absence of Nazi imagery, and multi-class classification: identifying specific symbols, while remaining
lightweight and adaptable enough to run on accessible hardware like single-GPU workstations. The same requirements
apply across adjacent domains, including online extremism tracking, digital forensics, and platform content
moderation—making the solution broadly relevant beyond its immediate scope.

\section{Objectives}

This research aims to design, implement, and evaluate an efficient and robust image classification framework for the
automatic detection and categorization of Nazi symbols in visual media. While the immediate focus is on Nazi
iconography, the broader goal is to develop a generalizable and extensible approach that can be adapted to other
categories of extremist or harmful visual symbols. This makes the framework more useful across different
contexts—including civil society monitoring, journalism, and platform moderation—where flexibility and scalability
are essential.

The specific objectives of this study are as follows:

\begin{itemize}
\item Construct a curated dataset of approximately 100,000 images, encompassing both Nazi-related and non-Nazi
imagery to support supervised training and evaluation.
\item Annotate and structure the dataset to support both binary (presence/absence) and multi-class (specific symbol)
classification tasks.
\item Investigate a diverse range of \ac{DL} and multimodal models—including \ac{CNN}s, \ac{ViT}, and OpenCLIP
embeddings—combined with classical \ac{ML} algorithms for downstream classification.
\item Evaluate model performance using standard classification metrics (e.g., accuracy, precision, recall, F1-score,
AUC), and analyze the comparative strengths and limitations of each approach.
\item Develop a lightweight, modular \ac{API} service to enable efficient model inference via HTTP requests, thereby
supporting real-world deployment and integration.
\end{itemize}

While this study is constrained in terms of data coverage and computational resources, it proposes a scalable
frameworkthat can be repurposed for other extremist symbol detection tasks by updating the dataset and retraining or
fine-tuning the models. This generalizable design makes it especially suitable for actors operating in fast-changing
threat landscapes or under limited infrastructure.

\section{Scope and Limitations}

The visual symbols associated with Nazi and neo-Nazi ideologies are numerous and varied, both historically and in
contemporary usage. Given this diversity, it is challenging to construct a comprehensive dataset that captures all
such symbols. Instead, this study focuses on a representative subset of the most visually distinctive and commonly
encountered symbols. Specifically, the following classes are included: \textit{Black Sun}, \textit{Flash and Circle},
\textit{Broken Sun Cross}, \textit{The Happy Merchant}, \textit{Hitler}, \textit{Nazi Salute},
\textit{Atomwaffen Division Emblem}, \textit{Blood \& Honor Emblem}, \textit{Celtic Cross}, \textit{Combat 18 Emblem},
\textit{Golden Dawn}, \textit{Hammerskins Emblem}, \textit{Identitarian Movement Emblem}, \textit{Kolovrat},
\textit{National Rebirth of Poland Emblem}, \textit{Popular Front Emblem}, \textit{Yellow Badge}, \textit{Siegrune},
\textit{\ac{SS} Skull}, \textit{Sturmabteilung Emblem}, \textit{Swastika}, and \textit{Wolfsangel}.

The classification system is designed to support two core tasks:

\begin{enumerate}
  \item \textbf{Binary Classification:} Determining whether a given image contains any Nazi-related symbol.
  \item \textbf{Multi-Class Classification:} Identifying which specific symbol (among 13 predefined categories) is
  present in a given image found to have Nazi-related symbology.
\end{enumerate}

In addition to dataset limitations, there were also hardware constraints. All experiments were conducted on a single
consumer-grade GPU (NVIDIA RTX 4060Ti) with 16 GB of video memory. This limited the ability to fine-tune large
models or run memory-intensive experiments—particularly with transformer-based architectures (e.g., ViT) or
multimodal embeddings (e.g., OpenCLIP). These models were used in this study, but only in lightweight or partially
frozen configurations to stay within the available hardware limits. It also restricted the feasibility
of training multiple models in parallel for ensembling. As a result, priority was given to computationally efficient
architectures and lightweight training procedures that could fit within the available memory and runtime budget.

Despite these constraints, this work seeks to demonstrate the feasibility of developing effective, accessible, and
ethically responsible tools for the detection of extremist visual content. It is intended to support researchers,
content moderators, and policymakers in their efforts to monitor and counteract the dissemination of hate imagery
online.

\section{Thesis Structure}

This thesis is organized as follows:

\begin{itemize}
  \item \textbf{\fullref{chapter:literature-review}} – Surveys existing research on Nazi symbol detection and
  classification, emphasizing methodological trends, key findings, and identifying research gaps addressed in this work.

  \item \textbf{\fullref{chapter:methodology}} – Describes the data collection
  process, dataset composition, preprocessing techniques, and the used classification algorithms. It also
  outlines the model training process, evaluation metrics, and justifies the chosen methodologies.

  \item \textbf{\fullref{chapter:ethical-considerations}} – Discusses the ethical implications of working with sensitive
  content, including concerns around data privacy, potential misuse, and broader responsibilities associated with
  automated hate symbol detection.

  \item \textbf{\fullref{chapter:results}} – Presents the outcomes of the experiments, including both
  quantitative and qualitative analyses of model performance for binary and multi-class classification tasks.

  \item \textbf{\fullref{chapter:discussion-and-analysis}} – Interprets the results in relation to the
  research objectives, compares findings to prior work, and discusses system limitations and potential improvements.

  \item \textbf{\fullref{chapter:model-deployment-and-api-service}} – Describes the deployment strategy
  for the best-performing models, including the architecture of the \ac{API} service, the \ac{CI/CD} pipeline, and
  details of the available endpoints.

  \item \textbf{\fullref{chapter:conclusion-and-future-work}} – Summarizes the thesis contributions, reflects on key
  insights gained, and outlines directions for future research in the field of Nazi symbol classification.
\end{itemize}
